\chapter{在线稀疏高斯过程回归}

\section{引言}

标准高斯过程回归在理论上具有完备的贝叶斯推断框架与良好的不确定性量化能力，然而其计算复杂度随训练数据量呈三次方增长，推断阶段的
存储与时间开销均难以满足在线学习场景的实时性需求。与此同时，移动机器人在连续运行过程中所积累的观测数据具有流式、增量的特点，
标准 GP 每次接收新数据后须对全量历史数据重新求解，这使其在实际部署中几乎不具备可扩展性。上述局限促使研究者转向稀疏近似与在线
推断方向，以在可接受的计算代价内保留 GP 的核心优势。

本章针对轨迹跟踪控制中残差动力学的在线建模需求，提出一种面向在线学习的稀疏高斯过程回归算法\textbf{OSVGP}。算法以变分推断为理论基础，通过诱导
点处的变分均值与协方差作为模型的充分统计量，将完整数据集的推断问题转化为低维变分参数的优化问题；在此基础上，设计支持增量更新的
变分参数递推机制，使模型能够随新观测数据持续修正后验估计而无需访问历史数据；进一步引入诱导点的增量扩充策略，在不重新训练的前提
下提升模型对新状态区域的表达能力。算法以系统名义模型的残差动力学为学习目标，在每一时刻同时输出点预测与预测不确定性，为第四章的
 GP-MPC 控制集成提供标准化接口。

\section{稀疏高斯过程及其变分推断}

标准高斯过程回归在训练集规模为 $N$ 时，需要对 $N \times N$ 的核矩阵进行 Cholesky
分解，时间复杂度为 $\mathcal{O}(N^3)$，存储复杂度为 $\mathcal{O}(N^2)$。这一瓶颈在
自动驾驶等实时系统中尤为突出——车辆运行过程中持续产生传感器数据，直接应用标准 GPR
难以满足在线更新的延时约束。稀疏高斯过程（Sparse Gaussian Process，SGP）通过引入一组
数量远少于训练集的\textbf{诱导点}（inducing points）来近似完整的后验分布，将计算复杂度
降低至可接受的范围。变分推断为稀疏高斯过程提供了严谨的理论基础。其核心思想是引入一个变分分
布 $q(u)$ 来近似诱导点函数值 $u$ 的真实后验，并通过最大化证据下界（ELBO）来优化该分布。
这一框架不仅将预测计算复杂度降至 $O(M^2)$（$M$为诱导点数量），还能够提供在线学习的可能性。

\subsection{稀疏高斯过程的两个条件分布}

设诱导输入集合 $Z = [z_1, z_2, \cdots, z_M]^\top \in \mathbb{R}^{M \times D}$，
$M \ll N$，对应的诱导变量为
\begin{equation}
    u = f(Z) = [f(z_1), f(z_2), \cdots, f(z_M)]^\top
\end{equation}
由高斯过程的定义，$(f(X), u)$，$(f(x^*), u)$ 服从联合高斯分布：
\begin{equation}
    \begin{pmatrix} f(X) \\ u \end{pmatrix} \sim \mathcal{N}
    \left(0, \begin{bmatrix} K_{ff} & K_{fu} \\ K_{uf} & K_{uu} \end{bmatrix}\right), \quad 
    \begin{pmatrix} f(x^*) \\ u \end{pmatrix} \sim \mathcal{N}
    \left(0, \begin{bmatrix} K_{**} & K_{*u} \\ K_{u*} & K_{uu} \end{bmatrix}\right)
\end{equation}
基于高斯过程的一致性，我们可以通过积分（边缘化）$u$联合分布$p(f^*, f, u)$恢复$p(f^*, f)$：
\begin{equation}
    p(f^*, f) = \int p(f^*, f, u) \mathrm{d}u, \quad p(u) = \mathcal N(0, K_{uu})
\end{equation}
这是精确的表达式，接下来我们介绍几乎所有稀疏方法采取的近似假设：\textbf{给定$u$的条件下$f$和$f^*$是条件独立的}。
\begin{equation}
    p(f^*, f) \approxeq q(^*, f) = \int q(f^* | u) q(f | u) p(u) \mathrm{d}u
\end{equation}
“诱导变量”这一名称的由来是：由于$f$和 $f^*$只能通过$u$进行信息传递，因此$u$起到了诱导训练案例与测试案例之间依赖关系的作用。


利用条件高斯公式，训练函数值和测试函数值在给定诱导变量 $u$ 条件下的精确分布为
\begin{equation} \label{eq:conditionals}
\begin{aligned}
    \text{训练条件（training conditional）}: &\quad p({f}|{u}) = \mathcal N \big (K_{fu} K_{uu}^{-1} u, 
    K_{ff} - Q_{ff} \big ) \\
    \text{测试条件（test conditional）}: &\quad p({f}_*|{u}) = \mathcal N \big (K_{*u} K_{uu}^{-1} u, 
    K_{**} - Q_{**} \big )
\end{aligned}
\end{equation}
其中 $Q_{ab} = K_{au} K_{uu}^{-1} K_{ub}$ 称为 Nyström 近似矩阵，刻画了由诱导点
所能解释的函数值方差。$K_{ff} - Q_{ff}$ 为残差协方差，反映了诱导点无法捕捉的
函数变化量。

大部分的稀疏方法都可以看作是对这两个条件分布的不同近似方法，我们接下来便引入DTC和FITC。

\subsection{DTC与FITC近似}

\textbf{DTC}\cite{csato2002sparse}基于$f = K_{fu} K_{uu}^{-1} u$这一映射，将两个条件分布定义为：
\begin{equation}
    q_{\mathrm{DTC}}(f|u)=\mathcal{N}(K_{fu}K_{uu}^{-1}{u},0),
    \quad q_{\mathrm{DTC}}({f_{*}}|{u})=p({f_{*}}|{u})
\end{equation}
DTC的近似之处在于其训练条件分布，给定$u$下，训练数据$f$是确定的，而其测试条件没有经过近似，是精确的。

其预测分布为
\begin{equation} \label{equ:DTC_pred}
    \begin{aligned}
    q_{\text{DTC}}(f_* | y) &= \mathcal N \left (
        Q_{*f} (Q_{ff} + \sigma^2 I)^{-1} y,
        K_{**} - Q_{*f} (Q_{ff} + \sigma^2 I)^{-1} Q_{f*}
        \right ) \\
        &= \mathcal N \left (
            \sigma^{-2} K_{*u} \Sigma K_{uf} y,
            K_{**} - Q_{**} + K_{*u} \Sigma K_{u, *}
            \right )
    \end{aligned}
\end{equation}
其中
\begin{equation}
    \Sigma = (\sigma^{-2} K_{uf} K_{fu} + K_{uu})^{-1}
\end{equation}
此处我们使用了逆矩阵公式，将需要取逆的矩阵的维度从$\mathbb R^{N \times N}$降为了$\mathbb R^{M \times M}$：
\begin{equation}
    \begin{aligned}
        (Q_{ff} + \sigma^2 I)^{-1} 
        &= (K_{{fu}} K_{{uu}}^{-1} K_{{uf}} + \sigma^2 I)^{-1} \\
        &\rightarrow \text{使用逆矩阵公式} \\
        &= \sigma^{-2} I - \sigma^{-4} K_{{fu}} (K_{{uu}} + \sigma^{-2} K_{{uf}} K_{{fu}})^{-1} K_{{uf}} \\
        &= \sigma^{-2} I - \sigma^{-4} K_{{fu}} \Sigma K_{{uf}} \\
        &\rightarrow \sigma^{-2} K_{{uf}} K_{{fu}} = \Sigma^{-1} - K_{{uu}} \\
        Q_{*f} (Q_{ff} + \sigma^2 I)^{-1}
        &= \sigma^{-2} K_{{*u}} K_{{uu}}^{-1} K_{{uf}} - \sigma^{-2} K_{{*u}} K_{{uu}}^{-1} (\Sigma^{-1} - K_{{uu}}) \Sigma K_{{uf}} \\
        &= \sigma^{-2} K_{{*u}} \Sigma K_{{uf}}
    \end{aligned}
\end{equation}

\textbf{FITC}\cite{snelson2007local}则假设给定$u$的条件下，$f_i$彼此之间相互独立
\begin{equation}
    p(y | f) \simeq q(y | u) = \mathcal N \left ( 
        K_{fu} K_{uu}^{-1} u, 
        \mathrm{diag} [K_{ff} - Q_{ff}] + \sigma^2 I
         \right )
\end{equation}
那么两个条件分布为：
\begin{equation}
    q_{\mathrm{FITC}}({f|u})
    = \prod_{i=1}^{n} p(f_i | u)
    =\mathcal{N}(K_{{fu}}K_{{uu}}^{-1}{u}, \mathrm{diag} [K_{ff} - Q_{ff}]),
    \quad q_{\mathrm{FITC}}({f_{*}}|{\bf u})=p({f_{*}}|{\bf u})
\end{equation}

FITC并非直接忽略方差，而是通过引入额外的独立性假设来近似训练条件分布。其先验形式为（这里我们暂时只考虑单个测试样本的情况）：
\begin{equation}
    q_{\text{FITC}}(f, f_*) = \mathcal N \left ( 0, \begin{bmatrix}
        Q_{ff} + \mathrm{diag} [K_{ff} - Q_{ff}] & Q_{f*} \\
        Q_{*f} & K_{**}
    \end{bmatrix} \right )
\end{equation}

最终我们可以得到其预测公式：
\begin{equation}
    \begin{aligned}
    q_{\text{FITC}}(f_* | y) &= \mathcal N \left (
        Q_{*f} (Q_{ff} + \Lambda)^{-1} y,
        K_{**} - Q_{*f} (Q_{ff} + \Lambda)^{-1} Q_{f*}
        \right ) \\
        &= \mathcal N \left (
            K_{*u} \Sigma K_{uf} \Lambda^{-1} y,
            K_{**} - Q_{**} + K_{*u} \Sigma K_{u, *}
            \right )
    \end{aligned}
\end{equation}
其中$\Sigma = (K_{uf} \Lambda^{-1} K_{fu} + K_{uu})^{-1}$， 
$\Lambda = \mathrm{diag} [K_{ff} - Q_{ff} + \sigma^2 I]$.

从DTC到FITC的演变，体现了稀疏高斯过程近似中对不确定性建模的深化——DTC以确定性映射为代价追求简洁，
FITC则通过条件独立假设保留了局部不确定性，从而在保持计算效率的同时提升了模型的泛化能力。

\subsection{高斯过程的变分推断}
上文基于直接的假设推导了DTC和FITC的预测公式，这缺乏一定的理论保证，因此我们从变分推断的角度，推导出DTC的
公式（FITC同理），并基于此我们可以设计在线更新算法。

设变分分布$q(u) = \mathcal N(m_u, S_u)$，变分后验为
\begin{equation}
    q(f, u) = p(f | u) q(u)
\end{equation}

在变分推断的框架下，我们需要最大化对数边缘似然$\log p(y)$，并由此引入证据下界和$q(u)$和$p(u)$之间的
KL散度。
\begin{equation}
    \begin{aligned}
        \log p(y) &= \log \iint p(y, f, u) df du \\
        &= \log \iint p(y, f, u) \frac{q(f, u)}{q(f, u)} df du \\
        &= \log \iint q(f, u) \frac{p(y, f, u)}{q(f, u)} df du \\
        &= \log \mathbb E_{q(f, u)} \left [ \frac{p(y, f, u)}{q(f, u)} \right ] \\
        &\geq \mathbb E_{q(f, u)} \left [ \log \frac{p(y, f, u)}{q(f, u)} \right ] \quad \text{(Jensen不等式)} \\
        &= \mathbb E_{q(f, u)} \left [ \log \frac{p(y|f, u)p(f, u)}{p(f|u)q(u)} \right ] \\
        &= \mathbb E_{q(f, u)} \left [ \log \frac{p(y|f)p(f|u)p(u)}{p(f|u)q(u)} \right ] \\
        &= \mathbb E_{q(f, u)} \left [ \log \frac{p(y|f)p(u)}{q(u)} \right ] \\
        &= \mathbb E_{q(f, u)} \left [ \log p(y|f) \right ] - \mathbb E_{q(f, u)} \left [ \log \frac{q(u)}{p(u)} \right ] \\
        &= \mathbb E_{q(f, u)} \left [ \log p(y|f) \right ] - \text{KL} \left [q(u) || p(u) \right ]
    \end{aligned}
\end{equation}
显然，最大化边缘似然就是最小化$q(u)$和$p(u)$之间的KL散度。

其中证据变分下界ELBO就是：
\begin{equation}
    \mathcal L(q) = \mathbb E_{q(f, u)} \left [ \log \frac{p(y, f, u)}{q(f, u)} \right ]
    = E_{q(f, u)} \left [ \log p(y|f) \right ] - \text{KL} \left [q(u) || p(u) \right ]
\end{equation}

我们考虑第一项
\begin{equation}
    \begin{aligned}
        \mathbb E_{q(f, u)} \left [ \log p(y|f) \right ] &= \int q(f, u) \log p(y|f) df du \\
        &= \int p(f|u) q(u) \log p(y|f) df du \\
        &= \int q(u) \int P(f|u) \log p(y|f) df du
    \end{aligned}
\end{equation}
其中
\begin{equation}
    \begin{aligned}
        y|f & \sim \mathcal N(f, \sigma^2) \\
        f|u & \sim \mathcal N(K_{fu}K_{uu}^{-1}u, K_{ff} - Q_{ff})
    \end{aligned}
\end{equation}
那么
\begin{equation}
    \log p(y|f) = -\frac{N}{2} \log 2\pi - N \log \sigma - \frac{1}{2\sigma^2} (y-f)^T(y-f)
\end{equation}
我们可以得到内部的积分
\begin{equation}
    \begin{aligned}
        \int p(f|u) \log p(y|f) df &= -\frac{N}{2} \log 2\pi - N \log \sigma - \frac{1}{2\sigma^2} \mathbb E_{p(f|u)} \left [ (y-f)^T(y-f) \right ] \\
        & \text{使用二次型的期望公式} \\
        &= -\frac{N}{2} \log 2\pi - N \log \sigma \\
        &- \frac{1}{2\sigma^2} \left [ (y - K_{fu}K_{uu}^{-1}u)^T(y - K_{fu}K_{uu}^{-1}u) + \text{Tr}(K_{ff} - Q_{ff}) \right ] \\
    \end{aligned}
\end{equation}
这样总的积分就是
\begin{equation}
    \begin{aligned}
        \mathbb E_{q(f, u)} \left [ \log p(y|f) \right ] &= -\frac{N}{2} \log 2\pi - N \log \sigma \\
        &- \frac{1}{2\sigma^2} \left [ (y - K_{fu}K_{uu}^{-1}m_u)^T(y - K_{fu}K_{uu}^{-1}m_u) + \text{Tr}(K_{ff} - Q_{ff}) \right ] \\
        &- \frac{1}{2\sigma^2} \text{Tr}(K_{fu}K_{uu}^{-1}S_u K_{uu}^{-1}K_{uf})
    \end{aligned}
\end{equation}

我们接着考虑第二项KL散度， $p(u) \sim \mathcal N(0, K_{uu})$， $q(u) \sim \mathcal N(m_u, S_u)$：
\begin{equation}
    \begin{aligned}
        -\text{KL} \left [q(u) || p(u) \right ] &=
        -\mathbb E_{q(f, u)} \left [ \log q(u) \right ] + \mathbb E_{q(f, u)} \left [ \log p(u) \right ] \\
        &= -\mathbb E_{q(u)} \left [ \log q(u) \right ] + \mathbb E_{q(u)} \left [ \log p(u) \right ]
    \end{aligned}
\end{equation}
其中
\begin{equation}
    \begin{aligned}
        \mathbb E_{q(u)} \left [ \log q(u) \right ] &= -\frac{M}{2} \log 2 \pi - \frac{1}{2} \log | S_u |
        -\frac{1}{2} \mathbb E_{q(u)} \left [ (u-m_u)^\top S_u^{-1} (u-m_u) \right ] \\
        &= -\frac{M}{2} \log 2 \pi - \frac{1}{2} \log | S_u | - \frac{M}{2}
    \end{aligned}  
\end{equation}
并且
\begin{equation}
    \begin{aligned}
        \mathbb E_{q(u)} \left [ \log p(u) \right ] &= -\frac{M}{2} \log 2 \pi - \frac{1}{2} \log | K_{uu} |
        -\frac{1}{2} \mathbb E_{q(u)} \left [ u^\top K_{uu}^{-1} u \right ] \\
        &= -\frac{M}{2} \log 2 \pi - \frac{1}{2} \log | K_{uu} | \\
        &-\frac{1}{2} \mathbb E_{q(u)} \left [ (m_u + (u-m_u))^\top K_{uu}^{-1} (m_u + (u-m_u)) \right ] \\
        &= -\frac{M}{2} \log 2 \pi - \frac{1}{2} \log | K_{uu} | - \frac{1}{2} m_u^T K_{uu}^{-1} m_u
        -\frac{1}{2} \text{Tr}(K_{uu}^{-1}S_u)
    \end{aligned}  
\end{equation}
那么
\begin{equation}
    \begin{aligned}
        -\text{KL} \left [q(u) || p(u) \right ] &=
        -\frac{1}{2} \left [ -\log |S_u| - M + \log |K_{uu}| + m_u^T K_{uu}^{-1} m_u + \text{Tr}(K_{uu}^{-1}S_u) \right ]
    \end{aligned}
\end{equation}

只保留关于$m_u$和$S_u$的项，ELBO的剩余的表达式是
\begin{equation}
    \begin{aligned}
        \tilde{\mathcal L(q)} &= \frac{1}{2} \left [
            \log |S_u| - \text{Tr}(K_{uu}^{-1}S_u) - m_u^T K_{uu}^{-1} m_u
        \right ] \\
        & -\frac{1}{2 \sigma^2} \left [
            -2y^\top K_{fu} K_{uu}^{-1} m_u + m_u^\top K_{uu}^{-1} K_{uf} K_{fu} K_{uu}^{-1} m_u
            + \text{Tr}(K_{fu}K_{uu}^{-1}S_u K_{uu}^{-1}K_{uf})
        \right ]
    \end{aligned}
\end{equation}

分别对$m_u$和$S_u$求导数：
\begin{equation}
    \begin{aligned}
        \frac{\partial \tilde{\mathcal L(q)}}{\partial m_u} &= 
        -K_{uu}^{-1} m_u + \sigma^{-2} K_{uu}^{-1} K_{uf} y - \sigma^{-2} 
        K_{uu}^{-1} K_{uf} K_{fu} K_{uu}^{-1} m_u \\
        \frac{\partial \tilde{\mathcal L(q)}}{\partial S_u} &= 
        \frac{1}{2} S_u^{-1} - \frac{1}{2} K_{uu}^{-1} - \frac{1}{2 \sigma^2} 
        K_{uu}^{-1} K_{uf} K_{fu} K_{uu}^{-1} \\
    \end{aligned}
\end{equation}
分别等于0得到最优值
\begin{equation}
    \begin{aligned}
        m_u &= \sigma^{-2} K_{uu} \Sigma K_{uf} y \\
        S_u &= K_{uu} \Sigma K_{uu}
    \end{aligned}
\end{equation}
其中
\begin{equation}
    \Sigma = (K_{uf} K_{fu} / \sigma^2 + K_{uu})^{-1}
\end{equation}

根据总期望公式和总方差公式我们可以得到：
\begin{equation}
    \begin{aligned}
    \mathbb E(f_* | y) &= \mathbb E_{q(u)} [\mathbb E(f_* | u)] = \mathbb E_{q(u)} [K_{*u} K_{uu}^{-1} u] = K_{*u} K_{uu}^{-1} m_u \\
    \text{Var}(f_* | y) &= \mathbb E_{q(u)} [\text{Var}(f_* | u)] + \text{Var}_{q(u)} [\mathbb E(f_* | u)] \\
    &= \mathbb E_{q(u)} [K_{**} - Q_{**}] + \text{Var}_{q(u)} [K_{*u} K_{uu}^{-1} u] \\
    &= K_{**} - Q_{**} + K_{*u} K_{uu}^{-1} S_u K_{uu}^{-1} K_{u*}
\end{aligned}
\end{equation}
我们便可以得到预测分布
\begin{equation} \label{equ:mu_su_predictive}
    \begin{aligned}
        q(f_* | y) &= \int p(f_* | u) q(u) du \\
        &= \mathcal N (K_{*u} K_{uu}^{-1} m_u, K_{**} - Q_{**} + K_{*u} K_{uu}^{-1} S_u K_{uu}^{-1} K_{u*}) \\
        &= \mathcal N (\sigma^{-2} K_{*u} \Sigma K_{uf} y, K_{**} - Q_{**} + K_{*u} \Sigma K_{u*})
    \end{aligned}
\end{equation}

这跟我们推导的DTC的\autoref{equ:DTC_pred}一模一样，但是我们得到了ELBO得到的理论下界保证。同理在FITC的框架下，
诱导变量的变分分布是
\begin{equation}
    \begin{aligned}
        m_u &= K_{uu} \Sigma K_{uf} \Lambda^{-1} y \\
        S_u &= K_{uu} \Sigma K_{uu}
    \end{aligned}
\end{equation}
其中
\begin{equation}
    \Sigma = (K_{uf} \Lambda^{-1} K_{fu} + K_{uu})^{-1}, \quad \Lambda = \text{diag}(K_{ff} - Q_{ff} + \sigma^2)
\end{equation}

\section{在线更新算法}

在许多实际应用场景中，数据通常以流式（streaming）的形式不断到达，传统的批处理（batch）推断方法需要反复在全量数据
上进行优化，计算代价较高且难以满足实时性需求。因此，有必要设计能够随数据逐步更新的在线推断算法。

本节围绕变分推断框架，讨论模型在新数据到达时的更新策略，根据是否需要扩充诱导变量及其分布分为非增量和增量更新算法。

\subsection{变分分布的非增量更新}

Huber等人\cite{Huber_2014}在标准GP模型上实现了一个类卡尔曼更新法，我们将这个方法推广至稀疏变分结构。

在$k-1$时刻，我们有变分分布$q(u) = \mathcal N (m_{k-1}, S_{k-1})$，在$k$时刻我们得到了一个新的数据点
$(x_k, y_k)$，$p(y_k | u) = \mathcal N(K_{f_k u} K_{uu}^{-1}u, \sigma^2)$，$u$在$y_k$下的条件概率（后验）为
\begin{equation}
    q(u|y_k) = \frac{q(y_k|u)q(u)}{p(y_k)} \propto q(y_k|u)q(u)
\end{equation}
取负对数
\begin{equation}
    \begin{aligned}
        - \log q(y_k|u)q(u)
        &= \frac{1}{2} (u - m_{k-1})^T S_{k-1}^{-1} (u - m_{k-1}) + \frac{1}{2\sigma^2} (y_k - K_{f_ku} K_{uu}^{-1} u)^2 + \text{常量} \\
        &\rightarrow \text{只保留依赖于} u \text{的项} \\
        &= \frac{1}{2} u^T (S_{k-1}^{-1} + \sigma^{-2} K_{u u}^{-1} K_{uf_k} K_{f_ku} K_{uu}^{-1}) u \\
        &- (S_{k-1}^{-1} m_{k-1} + \sigma^{-2} K_{uu}^{-1} K_{uf_k} y_k) u \\
        &= \frac{1}{2} u^T A u - b^T u + \text{常量} \\
        &= \frac{1}{2} (u - A^{-1} b)^T A (u - A^{-1} b) + \text{常量}
    \end{aligned}
\end{equation}
这个后验是一个指数二次型，这必然是一个高斯分布，于是我们得到后验分布$q(u|y_k) = \mathcal N(m_k, S_k)$：
\begin{equation}    
    \begin{aligned}
        S_k &= A^{-1} =  (S_{k-1}^{-1} + \sigma^{-2} K_{u u}^{-1} K_{u f_k} K_{f_k u} K_{u u}^{-1})^{-1} \\
        m_k &= A^{-1}b = S_k (S_{k-1}^{-1} m_{k-1} + \sigma^{-2} K_{u u}^{-1} K_{u f_k} y_k)
    \end{aligned}   
\end{equation}
我们可以使用逆矩阵定理处理这个逆：
\begin{equation}
    S_k = S_{k-1} - S_{k-1} K_{u u}^{-1} K_{u f_k} (\sigma^2 + K_{f_k u} K_{u u}^{-1} S_{k-1} K_{u u}^{-1} K_{u f_k})^{-1} K_{f_k u} K_{u u}^{-1} S_{k-1}
\end{equation}

这是一个贝叶斯更新步骤，我们可以将整个更新方案写成一个类似卡尔曼滤波的形式，定义$\phi_k = K_{f_k u} K_{uu}^{-1}$，那么
$y_k = \phi_k u + \epsilon, \quad \epsilon \sim \mathcal N(0, \sigma^2)$
\begin{equation}
    \begin{aligned}
        & r_k = y_k - \Phi_k m_{k-1} \\
        & s_k = \sigma^2 + \Phi_k S_{k-1} \Phi_k^T \\
        & K_k = S_{k-1} \Phi_k^T s_k^{-1} \\
        & m_k = m_{k-1} + K_k r_k \\
        & S_k = S_{k-1} - K_k \phi_k S_{k-1}
    \end{aligned}
\end{equation}
第一项是残差，第二项是预测方差，第三项是增益。总结一下即为\autoref{alg:online_vsgp_update}，其总体计算复杂度为$\mathcal O(M^2)$。

\begin{algorithm}[htbp]
    \caption{在线变分稀疏高斯过程更新}
    \label{alg:online_vsgp_update}
    \begin{algorithmic}[1]
        \Require 
        \Statex 上一时刻变分分布参数：均值 $m_{k-1}$，协方差 $S_{k-1}$
        \Statex 诱导点协方差矩阵 $K_{uu}$（固定）
        \Statex 新观测数据点 $(x_k, y_k)$
        \Statex 观测噪声方差 $\sigma^2$
        \Statex 核函数 $k(\cdot, \cdot)$
        
        \Function{OnlineUpdate}{$m_{k-1}$, $S_{k-1}$, $K_{uu}$, $x_k$, $y_k$, $\sigma^2$}
        
            \State \textbf{步骤1}: 计算核向量
            \State $K_{f_k u} \gets [k(x_k, u_1), k(x_k, u_2), \ldots, k(x_k, u_M)]$ \Comment{$M$为诱导点个数}
            
            \State \textbf{步骤2}: 计算线性化观测矩阵
            \State $\phi_k \gets K_{f_k u} K_{uu}^{-1}$ \Comment{将观测映射到诱导点空间}
            
            \State \textbf{步骤3}: 计算预测
            \State $\hat{y}_k \gets \phi_k m_{k-1}$ \Comment{预测观测值}
            \State $r_k \gets y_k - \hat{y}_k$ \Comment{残差（新息）}
            
            \State \textbf{步骤4}: 计算预测方差
            \State $s_k \gets \sigma^2 + \phi_k S_{k-1} \phi_k^\top$ \Comment{标量预测方差}
            
            \State \textbf{步骤5}: 计算增益
            \State $K_k \gets S_{k-1} \phi_k^\top s_k^{-1}$ \Comment{增益向量}
            
            \State \textbf{步骤6}: 更新后验分布
            \State $m_k \gets m_{k-1} + K_k r_k$ \Comment{更新均值}
            \State $S_k \gets S_{k-1} - K_k \phi_k S_{k-1}$ \Comment{更新协方差}
            
            \State \textbf{return} $m_k$, $S_k$
        \EndFunction
    \end{algorithmic}
\end{algorithm}

\subsection{变分分布的增量更新}

在标准的变分稀疏高斯过程回归中，若仅对变分分布进行非增量更新，则诱导点的位置保持不变，仅对其对应的变分均值与协方差进行调整。
在这种情况下，模型的表达能力受限于已有诱导点在输入空间中的分布范围。

当测试输入 $x^*$ 位于训练数据及诱导点未充分覆盖的区域时，由于缺乏有效的支撑信息，高斯过程的后验预测将逐渐退化为其先验分布，
即回归到先验均值 $m(x^*)$，同时伴随较大的不确定性。这种现象表明，固定诱导点位置的变分更新机制在面对输入分布扩展或新区域数据时，
难以提供准确的预测结果。

为了解决上述问题，有必要对变分分布进行增量扩展，在保留已有统计信息的基础上，引入新的诱导点以覆盖新的输入区域，从而提升模型在未
见区域的泛化能力与预测精度。

观察\autoref{equ:mu_su_predictive}我们发现，增加诱导变量会导致标准稀疏高斯过程预测分布的$K_{u u}$， $K_{u f}$这两项变化，
需要更新，从变分推断的角度则需要更新$K_{u u}$，$m_u$，$S_u$。

设新诱导点为$u_{n}$，原有诱导点为$u_{o}$，$K_{u f}$的更新比较简单：
\begin{equation}
    K_{u f} = \begin{bmatrix}
        K_{u_o f} \\
        K_{u_n f}
    \end{bmatrix}
\end{equation}

$K_{u u}$则并不会直接更新，而是更新对应的Cholesky分解矩阵$L_{u u}$，令
\begin{equation}
    C = L_{u_o u_o}^{-1} K_{u_o u_n}, \quad D = \sqrt{K_{u_n u_n} - C^\top C}
\end{equation}
则
\begin{equation}
    L_{u u} = \begin{bmatrix}
        L_{u_o u_o} & 0 \\
        C^\top & D
    \end{bmatrix}
\end{equation}

$m_u$的更新方法为
\begin{equation}
    m_u = \begin{bmatrix}
        m_{u_o} \\
        (L_{u_o u_o}^{-\top} L_{u_o u_o}^{-1} K_{u_o u_n})^\top m_{u_o}
    \end{bmatrix}
\end{equation}

$S_u$的更新方法为
\begin{equation}
    S_u = \begin{bmatrix}
        S_{u_o} & (S_{u_o} L_{u_o u_o}^{-\top} L_{u_o u_o}^{-1} K_{u_o u_n})^\top \\
        S_{u_o} L_{u_o u_o}^{-\top} L_{u_o u_o}^{-1} K_{u_o u_n} & S_{n n}
    \end{bmatrix}
\end{equation}
其中
\begin{equation}
    S_{n n} = K_{u_n u_n} - (L_{u_o u_o}^{-1} K_{u_o u_n})^\top (L_{u_o u_o}^{-1} K_{u_o u_n}) + 
    (L_{u_o u_o}^{-\top} L_{u_o u_o}^{-1} K_{u_o u_n})^\top S_{u_o} (L_{u_o u_o}^{-\top} L_{u_o u_o}^{-1} K_{u_o u_n})
\end{equation}

这是由$u_n = K_{u_n u_o} K_{u_o u_o}^{-1} S_{u_o}$得来的。

但我们还需要知道何时加入一个新的诱导点，一个直觉上的做法是避免空间冗余，如果新诱导点跟旧诱导点在空间上相距较远则加入。
但这种在欧氏空间下的距离度量并不是稀疏高斯所关心的，稀疏高斯过程利用诱导点构建了一个低维子空间$\mathcal H_U = \text{span}
\{ k(\cdot, u_i) \}$，关键是新点是否能够提供新的基方向，是在函数空间（再生核希尔伯特空间，RKHS）里的新方向。

为此我们需要计算这种相关性，计算当前诱导点子空间解释了新诱导点多少比例的方差，即
\begin{equation}
    r(u_o) = \frac{K_{u_n u_n} - K_{u_n u_o} K_{u_o u_o}^{-1} K_{u_o u_n}}{K_{u_n u_n}}
\end{equation}
当该值大于给定的阈值$\epsilon$时，就加入新诱导点。

我们将上述增量更新方法记为\autoref{alg:add_new_inducing}，其计算复杂度仍为$\mathcal O(M^2)$ （变分法不需要扩充$k_{u f}$）。

结合两个算法我们可以得出如下结论：OSVGP的计算复杂度为$\mathcal O(M^2)$，无论是增量还是非增量方案。

\begin{algorithm}[htbp]
    \caption{诱导点增量更新}
    \label{alg:add_new_inducing}
    \begin{algorithmic}[1]
        \Require 
        \Statex 旧诱导变量$u_o$及其协方差矩阵的分解矩阵$L_{u_o u_o}$，阈值参数$\epsilon$;
        新旧诱导点间的协方差$K_{u_o u_n}$，新诱导点自身的协方差$K_{u_n u_n}$;
        旧变分分布参数：均值 $m_{u_o}$，协方差 $S_{u_o}$

        \Function{CheckRelevance}{$L_{u_o u_o}$, $K_{u_o u_n}, \epsilon$}
            \State $C \gets L_{u_o u_o}^{-1} K_{u_o u_n}$
            \State $v \gets K_{u_n u_n} - C^\top C$
            \State $r(u_n) \gets \frac{v}{K_{u_n u_n}}$
            \State \textbf{return} $r(u_n) > \epsilon$
        \EndFunction

        \Function{AddNewInducing}{$u_n$, $L_{u_o u_o}$, $K_{u_o u_n}$, $K_{u_n u_n}$, $m_{u_o}$, $S_{u_o}$}
            
            \State \textbf{步骤1}: 计算中间矩阵C和D
            \State $C \gets L_{u_o u_o}^{-1} K_{u_o u_n}$
            \State $D \gets \sqrt{K_{u_n u_n} - C^\top C}$
            
            \State \textbf{步骤2}: 更新Cholesky因子$L_{uu}$
            \State $L_{uu} \gets \begin{bmatrix} L_{u_o u_o} & 0 \\ C^\top & D \end{bmatrix}$
            
            \State \textbf{步骤3}: 更新均值$m_u$
            \State $T \gets L_{u_o u_o}^{-\top} L_{u_o u_o}^{-1} K_{u_o u_n}$ \Comment{临时变量}
            \State $m_u \gets \begin{bmatrix} m_{u_o} \\ T^\top m_{u_o} \end{bmatrix}$
            
            \State \textbf{步骤4}: 更新协方差$S_u$的右下块$S_{nn}$
            \State $A \gets L_{u_o u_o}^{-1} K_{u_o u_n}$ \Comment{临时变量}
            \State $S_{nn} \gets K_{u_n u_n} - A^\top A + T^\top S_{u_o} T$
            
            \State \textbf{步骤5}: 组装完整的$S_u$矩阵
            \State $S_{u_o u_n} \gets S_{u_o} T$
            \State $S_u \gets \begin{bmatrix} S_{u_o} & S_{u_o u_n}^\top \\ S_{u_o u_n} & S_{nn} \end{bmatrix}$
            
            \State \textbf{return} $L_{uu}$, $m_u$, $S_u$
        \EndFunction
    \end{algorithmic}
\end{algorithm}

\section{实验数据和分析}

\subsection{数据集选择与评估指标}

本章的数据有两个来源，一个是用于演示算法效果，对比标准高斯过程和其他在线算法的仿真数据，另一个是在实车测试平台上获取的多维数据。
前者我们采用以下函数生成数据：
\begin{equation}
    y = x \sin x + w, \quad w \sim \mathcal N(0, 0.33^2)
\end{equation}
而后者则同第四章 \autoref{equ:augmented_zk} 所示
\begin{equation}
    z_k = [v_k,\ w_k,\ a_k,\ \delta_k,\ v_{k-1},\ w_{k-1},\ a_{k-1},\ \delta_{k-1},\ \ldots,\ 
    v_{k-T+1},\ w_{k-T+1},\ a_{k-T+1},\ \delta_{k-T+1}]^\top
\end{equation}
这是一个$D=4T$维的特征量，对应系统残差$d_k$为需预测值$y_k$, 那么
\begin{equation}
X = [z_1, z_2, \ldots, z_N]^\top \in \mathbb R^{N \times D}, 
\quad Y = [y_1, y_2, \ldots, y_N]^\top \in \mathbb R^{N \times 1}
\end{equation}

为了全面评估高斯过程回归模型的预测性能，本文采用三种互补的评价指标：均方根误差（Root Mean Squared Error, RMSE）、平均绝对
误差（Mean Absolute Error, MAE）以及负对数似然（Negative Log-Likelihood, NLL）。这些指标分别从点预测精度和不确定性建模质量
两个方面进行评价。

\textbf{均方根误差（RMSE）}用于衡量预测均值与真实值之间的偏差，其定义为：
\begin{equation}
\mathrm{RMSE} = \sqrt{ \frac{1}{n} \sum_{i=1}^{n} (y_i - \mu_i)^2 },
\end{equation}
其中，$y_i$ 表示真实值，$\mu_i$ 表示模型的预测均值。RMSE对较大的预测误差更为敏感。

\textbf{平均绝对误差（MAE）}用于衡量预测值与真实值之间的平均绝对偏差，其定义为：
\begin{equation}
\mathrm{MAE} = \frac{1}{n} \sum_{i=1}^{n} |y_i - \mu_i|,
\end{equation}
其中，$y_i$ 和 $\mu_i$ 分别表示真实值和预测均值。相比RMSE，MAE对异常值更加鲁棒。

\textbf{负对数似然（NLL）}，高斯过程模型输出预测分布 $y_i \sim \mathcal{N}(\mu_i, \sigma_i^2)$，NLL用于评估预测分布对真实数
据的拟合程度，其定义为：
\begin{equation}
\mathrm{NLL} = \frac{1}{n} \sum_{i=1}^{n} \left[ 
\frac{1}{2} \log (2\pi \sigma_i^2) + \frac{(y_i - \mu_i)^2}{2\sigma_i^2}
\right],
\end{equation}
其中 $\sigma_i^2$ 为预测方差。NLL同时考虑预测误差与不确定性估计的合理性，值越小表示模型的概率预测性能越好。

综上，RMSE和MAE用于评估点预测精度，其中RMSE对大误差更敏感，而MAE具有更强的鲁棒性；NLL则进一步评估模型预测分布的质量，从而综合反映模型
的预测性能与不确定性建模能力。

除此以外我们还需要对比不同算法的计算时间，对比这一指标对实时系统上运行的算法非常关键。本章的实验是在个人PC上执行Python代码得到的结果，
处理器为英特尔 i5-13400，内存32GB。

\subsection{稀疏化与变分推断效果}

给定平方指数核
\begin{equation}
    k(x,x') =
    0.33^2
    \exp
    \left(
    -\frac{(x-x')^2}{2}
    \right)
\end{equation}
我们实现标准高斯过程回归，稀疏高斯过程（DTC 和 FITC），以及 FITC 的变分推断模型，
在合成数据集上进行对比实验。训练集由函数 $y = x\sin x$ 加噪声生成，共 51 个样本点，
噪声标准差为 $\sigma = 0.33$；诱导点从训练集中均匀抽取 11 个。

如~\autoref{fig:compare_4gp} 所示，四种方法均能在训练区间 $[0, 20]$ 内捕捉到
函数的主要趋势。标准 GP（左上）的均值曲线与数据点吻合最为紧密，置信区间在训练区间
内较窄，而在外推区域 $(-5, 0)$ 和 $(20, 25)$ 迅速展宽，体现出合理的不确定性估计。
DTC GP（右上）和 FITC GP（左下）均使用相同的 11 个诱导点，均值曲线整体走势与标准
GP 相近，但置信区间在训练区间内略宽，说明稀疏近似引入了额外的近似误差。
FITC 变分推断模型（右下）的预测结果与 FITC GP 一致，表明在当前实验设置下两者
收敛至相同的解。

\begin{figure}[htbp]
    \centering
    \includegraphics[width=1\linewidth]{onlinegp/compare_4gp.png}
    \caption{不同 GP 模型的预测结果对比。红色曲线为预测均值，红色阴影为 $\pm 2\sigma$
    置信区间，蓝色圆点为诱导点，叉号为训练数据。}
    \label{fig:compare_4gp}
\end{figure}

\autoref{tab:gp_metrics} 列出了四种模型在训练集上的定量评估结果。

\begin{table}[htbp]
    \centering
    \caption{不同 GP 模型的训练集误差指标}
    \label{tab:gp_metrics}
    \begin{tabular}{lcccc}
        \toprule
        模型 & RMSE & MAE & NLL & Time (ms) \\
        \midrule
        标准 GP        & 0.424 & 0.316 & 1.240 & 1.058 \\
        DTC GP         & 1.403 & 1.079 & 6.448 & 0.643 \\
        FITC GP        & 1.587 & 1.172 & 5.396 & 0.618 \\
        变分 FITC GP   & 1.587 & 1.172 & 5.396 & 0.599 \\
        \bottomrule
    \end{tabular}
\end{table}

由表可见，标准 GP 在三项指标上均显著优于两种稀疏方法，RMSE 仅为 0.424，NLL 为
1.240。稀疏方法的误差较大，主要原因在于诱导点数量（11 个）相对于训练集规模（51
个）较少，信息压缩导致近似精度下降。DTC 与 FITC 的 NLL 分别为 6.448 和 5.396，
FITC 略优，这与 FITC 对角修正项能更好地保留训练点处的预测不确定性的理论分析一致。
变分 FITC GP 与 FITC GP 结果完全相同，说明在固定超参数、不进行变分优化的条件下，
两者退化为等价模型。

最后一列的计算时间是执行一百次取平均值得到的结果，可以看到稀疏方法在速度上均比标准高斯实现有优势。以上结果均
符合本章的理论分析。

\begin{figure}[htbp]
    \centering
    \includegraphics[width=1\linewidth]{onlinegp/trained.png}
    \caption{DTC法超参数和诱导点未训练（左）和训练（右）的预测图}
    \label{fig:trained}
\end{figure}

\autoref{fig:trained} 进一步对比了DTC GP在超参数与诱导点位置优化前后的预测效果，
定量结果见\autoref{tab:trained_metrics}。

\begin{table}[htbp]
    \centering
    \caption{DTC GP 训练前后的误差指标对比}
    \label{tab:trained_metrics}
    \begin{tabular}{lccc}
        \toprule
        模型 & RMSE & MAE & NLL \\
        \midrule
        未训练 & 1.404 & 1.080 & 6.449 \\
        训练后 & 0.267 & 0.211 & 0.201 \\
        \bottomrule
    \end{tabular}
\end{table}

经过超参数优化与诱导点位置联合训练后，模型性能得到大幅提升：RMSE 由 1.404 降至
0.267，降幅达 81.0\%；NLL 由 6.449 骤降至 0.201，接近标准 GP 的水平（1.240），
说明训练后的稀疏模型在保持计算效率优势的同时，预测精度与不确定性估计质量均得到
显著改善。

\subsection{在线学习算法效果对比}

本小节对四种高斯过程在线学习方法进行系统性对比实验，包括：标准非在线高斯过程（Standard GP，作为基线）、本文提出的OSVGP算法、Bijl\cite{bijl_2015_online}提出的
基于在线FITC/PITC近似的序贯更新算法，以及Gómez-Verdejo\cite{gomez_asgp}等人提出的自适应高斯过程（AGP）算法。
选取上述两种对比方法的原因如下：Bijl方法基于FITC/PITC稀疏近似框架，通过序贯贝叶斯更新实现在线学习，其推断路线与本文基于变分推断（ELBO）的OSVGP有本质区别，可
作为非变分在线稀疏GP的代表性基线；AGP方法虽源于变分稀疏GP框架，但本文所采用的为其不更新超参数的快速版本，每次新样本到来时仅自适应更新诱导点集合而不执行推断
迭代，计算开销极低，代表了另一类轻量级在线学习策略。
更重要的是，两种算法均能在实际部署中满足本系统20\,Hz的实时运行频率需求，因此具有直接的可比性。
实验旨在从预测精度与计算效率两个维度评估OSVGP算法的在线学习性能。所有计时结果均为执行一百次取平均值。
值得注意的是，本文中AGP的非增量更新算法并非真正的非增量更新，而只是说诱导点是非增量的，其算法还是会增量更新历史数据。
AGP算法的遗忘参数~$\lambda$~设置为~$0.99$，诱导点扩充阈值参数~$\epsilon = 0.01$。

\subsubsection{同域非增量更新}

第一组实验考察各算法在原始数据域内的非增量更新性能。训练数据与更新数据均采样自 $x \in [0, 20]$，更新数据对应的目标函数退化
为 $0.33 \cdot x\sin(x) + w$，以模拟系统动态缓慢漂移的场景。实验结果如\autoref{fig:compare_in_same} 所示，定量指标
见\autoref{tab:same_domain}。

\begin{figure}[htbp]
    \centering
    \includegraphics[width=1\linewidth]{onlinegp/compare_in_same.png}
    \caption{不同在线GP模型在原数据域执行非增量更新的预测效果}
    \label{fig:compare_in_same}
\end{figure}

\begin{table}[htbp]
\centering
\caption{同域非增量更新实验结果}
\label{tab:same_domain}
\begin{tabular}{lcccc}
\toprule
方法 & RMSE $\downarrow$ & MAE $\downarrow$ & NLL $\downarrow$ & 时间 (ms) \\
\midrule
Standard GP & 4.5299 & 3.5319 & 101.3939 & —      \\
OSVGP       & 2.3181 & 1.8053 & 46.9220  & \textbf{6.958}  \\
Bijl        & 2.8124 & 2.1833 & 68.0426  & 7.682  \\
AGP         & 1.1860 & 0.9857 & 4.2223   & 8.142  \\
\bottomrule
\end{tabular}
\end{table}

在同域更新场景下，AGP算法凭借其滑动窗口遗忘机制，在三项指标上均取得压倒性优势，对函数缓慢漂移的适应
能力最强，这是因为AGP还是会增量更新（即便诱导点没有增量更新）和遗忘历史数据，由于新旧数据数量一致，使得AGP算法能够完整替换
所有历史数据，严格来说并不是非增量更新。

OSVGP相较于Standard GP基线在RMSE、MAE和NLL上分别提升了48.8\%、48.9\%和53.7\%，并以6.958~ms的耗时在三种在线方法中运行
最快。Bijl方法的预测精度虽优于Standard GP，但整体落后于OSVGP，说明本文所提的稀疏变分在线更新策略在同域场景下具有更强的后验修正能力。

\subsubsection{跨域非增量更新}

第二组实验模拟更具挑战性的跨域更新场景：训练数据来自 $x \in [0, 10]$，更新数据扩展至新区域 $x \in [5, 20]$，目标函数同时变化
为 $0.3 \cdot x\sin(x) + w$。此设置旨在评估各算法在输入空间发生外推时的泛化能力与更新有效性。实验结果如\autoref{fig:compare_no_same} 所
示，定量指标见\autoref{tab:diff_domain}。

\begin{figure}[htbp]
    \centering
    \includegraphics[width=1\linewidth]{onlinegp/compare_no_same.png}
    \caption{不同在线GP模型在新数据域执行非增量更新的预测效果}
    \label{fig:compare_no_same}
\end{figure}

\begin{table}[htbp]
\centering
\caption{跨域非增量更新实验结果}
\label{tab:diff_domain}
\begin{tabular}{lcccc}
\toprule
方法 & RMSE $\downarrow$ & MAE $\downarrow$ & NLL $\downarrow$ & 时间 (ms) \\
\midrule
Standard GP & 2.9502 & 2.4940 & 16.1597 & —      \\
OSVGP       & \textbf{2.5824} & \textbf{2.0846} & \textbf{8.2391}  & \textbf{7.216}  \\
Bijl        & 2.6267 & 2.1486 & 10.5021 & 7.995  \\
AGP         & 4.7083 & 3.8327 & 12.0032 & 8.584  \\
\bottomrule
\end{tabular}
\end{table}

在跨域更新场景下，结果呈现出与同域场景截然不同的规律。OSVGP在RMSE、MAE和NLL三项指标上均取得最优，说明其后验更新机制能够有效利用新域观
测数据修正诱导点处的后验分布，在输入空间发生外推时仍具备良好的泛化能力。AGP的表现则出现明显退化，RMSE（4.708）甚至高于Standard GP基线
（2.950），这是由于其遗忘机制在跨域场景下对历史知识的过度丢弃，使模型无法保留对原始数据域的有效表达。Bijl方法的各项指标介于OSVGP与
Standard GP之间，整体稳定但缺乏竞争力。上述结果表明，OSVGP在处理数据域发生漂移的跨域场景时，相较于依赖遗忘机制的AGP具有更强的鲁棒性。

此外，观察\autoref{fig:compare_no_same} 可以发现，所有方法在 $x > 20$ 的区域预测均逐渐回归至先验均值，
置信区间也随之扩张。这一现象是稀疏高斯过程的固有性质：在缺乏诱导点覆盖的输入区域，
模型无法从数据中获取有效的后验信息，预测自然退化为先验分布。
这也从侧面说明了诱导点布局对稀疏GP预测质量的关键影响，
并引出了第三组实验中动态扩充诱导集合的必要性。

\subsubsection{跨域增量更新}

第三组实验在第二组的基础上，进一步引入\textbf{诱导点增量扩充}机制：
在线更新过程中，当新数据点处的预测不确定性超过先验方差的设定阈值时，
将其动态纳入诱导集合，以增强模型对新域的表达能力。
实验结果如\autoref{fig:compare_online} 所示，定量指标见\autoref{tab:incremental}。

\begin{figure}[htbp]
    \centering
    \includegraphics[width=1\linewidth]{onlinegp/compare_online.png}
    \caption{不同在线GP模型在新数据域执行增量与非增量更新的预测效果对比}
    \label{fig:compare_online}
\end{figure}

\begin{table}[htbp]
\centering
\caption{跨域增量更新实验结果}
\label{tab:incremental}
\begin{tabular}{lcccc}
\toprule
方法 & RMSE $\downarrow$ & MAE $\downarrow$ & NLL $\downarrow$ & 时间 (ms) \\
\midrule
Standard GP & 2.9502 & 2.4940 & 16.1597 & —      \\
OSVGP       & \textbf{0.5716} & \textbf{0.3896} & 6.3140  & \textbf{27.032} \\
Bijl        & 1.4022 & 1.2262 & 36.7192 & 28.346 \\
AGP         & 1.7826 & 1.1257 & \textbf{2.6709}  & 27.785 \\
\bottomrule
\end{tabular}
\end{table}

引入诱导点动态扩充后，各方法的预测精度均较非增量场景有所提升，
计算耗时相应上升至约27--28~ms，主要来源于诱导集合规模增大所带来的矩阵运算量增加。

从\autoref{fig:compare_online} 可以直观地看出各方法的差异。
OSVGP的诱导点随在线更新过程沿新数据域自适应密集分布，
均值曲线对更新数据拟合效果最佳，置信区间紧致合理，
这也与其在RMSE（0.572）和MAE（0.390）上的最优表现相吻合。
Bijl方法同样完成了诱导点的扩充，但在输入空间右侧区域（$x > 15$）
均值曲线出现明显偏差，置信区间偏大，反映出其后验协方差在诱导点
快速扩充条件下难以得到准确维护，NLL（36.719）的显著劣化亦印证了这一点。
AGP的诱导点分布较为分散，均值曲线在中间过渡区域拟合质量欠佳；
其NLL（2.671）虽为三种在线方法中最优，但这在一定程度上源于
其置信区间较宽、对数据点的覆盖更为保守，而非点预测精度的真实提升，
RMSE（1.783）与MAE（1.126）均逊于OSVGP亦佐证了这一判断。

\subsubsection{小结}

综合三组实验可以得出以下结论：在同域非增量更新场景中，OSVGP的预测精度仅次于AGP，但计算耗时最低；在跨域非增量更新场景中，OSVGP在全部三项指
标上均取得最优，展现出对数据域漂移的强鲁棒性；在引入诱导点增量扩充的跨域场景中，OSVGP的预测精度获得大幅提升，综合表现最优。需要指出的是，本
节实验所用数据规模较小，各方法间的绝对时间差距尚不明显；随着数据维度与规模的增大，基于稀疏变分框架的OSVGP算法在每步更新的计算复杂度上具有明
显优势，其速度收益将随数据规模的增长而愈加显著。后续将在真实多维传感器数据上进一步验证上述结论。

\subsection{实车数据测试}

本小节在真实车辆采集的多维数据上对三种在线更新算法进行验证。特征量设置为12维
（即预测时域 $T=3$ 的情况），输入特征涵盖车辆状态量与控制输入的历史序列。
实验流程如下：首先利用车辆正常运行工况下采集的数据训练初始GP模型（超参数未经优化），
随后在两组异常工况数据上执行在线自适应更新，以评估各算法在真实场景下的适应能力。
两组测试工况分别为：（1）\textbf{50ms通信延迟工况}，模拟控制回路中存在固定延迟时
残差分布发生偏移的场景；（2）\textbf{单侧车轮驱动力不足工况}，模拟车辆执行器发生
部分失效时系统动态特性发生显著变化的场景。

训练数据集包含234个样本点及40个诱导点，数据维度如下：
\begin{equation}
    X_o \in \mathbb{R}^{234 \times 12}, \quad Y_o \in \mathbb{R}^{234 \times 1}, 
    \quad U_o \in \mathbb{R}^{40 \times 12}
\end{equation}
两组测试数据分别包含263和242个数据点，实验将测试数据逐点顺序喂入在线更新算法，
以严格模拟实时在线学习的应用场景。

\autoref{fig:bad_wheel_delay} 展示了初始模型（更新前）在两组测试工况上的预测效果。
可以观察到，未经更新的模型在两种异常工况下均存在无法执行预测，完全回归了先验，均值预测为0，方差预测很大。

\begin{figure}[htbp]
    \centering
    \includegraphics[width=0.75\linewidth]{onlinegp/bad_wheel_delay.png}
    \caption{未更新模型在两组工况数据上的预测效果。上图为单侧车轮驱动力不足工况，
    下图为50ms通信延迟工况。}
    \label{fig:bad_wheel_delay}
\end{figure}

\autoref{fig:error_loc} 展示了三种算法在在线更新过程中预测误差随数据点序号的变化
曲线，图中星形标记点为触发诱导点扩充条件的位置。可以看出，各算法的预测误差总体
上随更新数据的逐步接入而呈下降趋势（除了AGP算法在工况一下出现预测失常），且诱导点扩充事件主要集中于更新初期，
反映了模型在接触新工况数据初期不确定性较高、需要扩充表达能力的规律。

\begin{figure}[htbp]
    \centering
    \begin{subfigure}{0.475\textwidth}
        \centering
        \includegraphics[width=\linewidth]{onlinegp/bad_wheel_loc.png}
        \caption{单侧车轮驱动力不足工况}
        \label{fig:bad_wheel_loc}
    \end{subfigure}
    \hfill
    \begin{subfigure}{0.475\textwidth}
        \centering
        \includegraphics[width=\linewidth]{onlinegp/delay_loc.png}
        \caption{50ms通信延迟工况}
        \label{fig:delay_loc}
    \end{subfigure}
    \caption{在线更新过程中各算法预测误差随数据序号的变化曲线，
    星形标记点为触发诱导点动态扩充的位置。}
    \label{fig:error_loc}
\end{figure}

\autoref{fig:update_all} 及\autoref{tab:real_car} 给出了三种算法完成全部在线更新后
的最终预测效果与定量指标。

\begin{figure}[htbp]
    \centering
    \includegraphics[width=1.0\linewidth]{onlinegp/update_all.png}
    \caption{三种在线GP模型在两种工况数据下完成更新后的预测效果对比。}
    \label{fig:update_all}
\end{figure}

\begin{table}[htbp]
\centering
\caption{实车数据在线更新实验结果}
\label{tab:real_car}
\begin{tabular}{llcccc}
\toprule
工况 & 方法 & RMSE $\downarrow$ & MAE $\downarrow$ & NLL $\downarrow$ & 时间 (ms) \\
\midrule
\multirow{3}{*}{延迟工况}
 & OSVGP & \textbf{0.03899} & \textbf{0.02837} & \textbf{-1.5689} & \textbf{129.3} \\
 & Bijl  & 0.04125          & 0.03086          & -1.5500          & 199.5          \\
 & AGP   & 0.93364          & 0.41083          & 38.0038          & 156.7          \\
\midrule
\multirow{3}{*}{驱动力不足工况}
 & OSVGP & \textbf{0.03555} & \textbf{0.02759} & \textbf{-1.5646} & \textbf{118.7} \\
 & Bijl  & 0.04530          & 0.03766          & -1.4888          & 193.7          \\
 & AGP   & 0.03905          & 0.02994          & -1.0567          & 134.3          \\
\bottomrule
\end{tabular}
\end{table}

从实验结果可以得出以下结论。

\textbf{延迟工况：}OSVGP在RMSE、MAE和NLL三项指标上均优于Bijl方法，
同时计算耗时（129.3~ms）较Bijl（199.5~ms）降低约35\%。AGP在本工况下
出现严重退化（RMSE~=~0.934，NLL~=~38.00），与前述跨域实验中的分析一致，
其滑动窗口遗忘机制在数据分布发生较大偏移时容易导致历史知识的过度丢弃，
使模型在有限新数据条件下难以重建有效的后验表达。

\textbf{驱动力不足工况：}三种方法均取得较好的更新效果，OSVGP仍以最低的
RMSE（0.0356）和MAE（0.0276）位居首位，NLL同样最优（-1.565），且耗时
最短（118.7~ms）。Bijl方法的精度略有下降，AGP在本工况下表现较延迟工况
有所恢复，但NLL（-1.057）仍逊于OSVGP与Bijl，说明其置信区间估计的
可靠性相对较弱。

综合两组实车工况实验，OSVGP在预测精度与计算效率两个维度上均表现最优，
验证了本文所提在线更新算法在真实高维非线性系统中的有效性与工程实用性。
值得指出的是，本实验中超参数未经专门优化，若进一步对核函数参数进行离线
辨识，预测精度有望得到进一步提升。

\subsection{总结}
本章围绕高斯过程在线学习算法，从仿真数据与实车数据两个层面开展了系统性实验验证。
首先，通过稀疏化与变分推断的对比实验表明，标准GP的预测精度最高，但稀疏方法在经
过超参数与诱导点联合优化后，可在保持计算效率优势的同时大幅缩小与标准GP的性能差距，验证了稀疏近似框架的有效性。
其次，在三组在线学习对比实验中，本文提出的OSVGP算法展现出综合性能优势。在同域非增量更新场景下，OSVGP以最低计算耗时取得仅次于AGP的预测精
度；在跨域非增量更新场景下，OSVGP在全部三项指标上均优于其他方法，表现出对数据域漂移的强鲁棒性；在引入诱导点动态扩充的跨域增量更新场景下，
OSVGP的RMSE和MAE进一步大幅提升，综合表现最优。
最后，在真实车辆采集的12维传感器数据上，OSVGP在延迟工况与驱动力不足工况两组测试中均取得最优的预测精度与不确定性估计质量，同时计算耗时
相较Bijl方法和AGP算法均较低，验证了该算法在高维非线性实时系统中的有效性与工程实用性。